\documentclass[11pt]{article}
\usepackage[margin=1in]{geometry}
\usepackage{amsmath,amssymb,amsthm}
\usepackage{microtype}
\usepackage{enumitem}
\setlist{nosep,leftmargin=*}
\usepackage{booktabs}
\usepackage{xurl}
\usepackage[colorlinks=true,allcolors=blue]{hyperref}
\sloppy
\newtheorem{theorem}{Theorem}
\newtheorem{lemma}[theorem]{Lemma}
\newtheorem{proposition}[theorem]{Proposition}
\newtheorem{corollary}[theorem]{Corollary}
\theoremstyle{definition}
\newtheorem{definition}[theorem]{Definition}
\theoremstyle{remark}
\newtheorem{remark}[theorem]{Remark}
\newcommand{\F}{\mathbb{F}}
\newcommand{\Z}{\mathbb{Z}}
\newcommand{\Hom}{\operatorname{Hom}}
\newcommand{\Col}{\operatorname{Col}}
\newcommand{\cN}{\mathcal{N}}
\newcommand{\sharpP}{\#\mathrm{P}}
\newcommand{\FP}{\mathrm{FP}}
\setlength{\parskip}{0.3em}
\title{Counting homomorphisms from knot groups\\ to solvable groups can be $\sharpP$-hard}
\author{John Cullen}
\date{6 November 2026}
\begin{document}
\maketitle
\begingroup\renewcommand{\thefootnote}{}\footnotetext{This note was prepared
with the assistance of AI systems: Anthropic's Claude for the proofs, code and
text, and OpenAI's Codex for two independent reviews of the underlying proofs.
Every cited computation can be rerun from the public repository
\url{https://github.com/jscul/knotlab}; Section~\ref{sec:evidence} explains how.}\endgroup

\begin{abstract}
For a finite group $G$ and a link $L\subset S^3$, let $H_L(G)$ be the number of
homomorphisms from $\pi_1(S^3\setminus L)$ to $G$. Rowell conjectured in 2009
that $H_L(G)$ can be computed exactly in polynomial time from a diagram of $L$
whenever $G$ is solvable. Unless every function in $\sharpP$ is computable in
polynomial time, this fails. On links it fails already for $S_3$ and the other
dihedral groups of order $2n$ with $n$ odd: tying torus knots into the
components and cabling some of them turns the count into a Potts model, at the
price of unboundedly many components. The main point of this note is that it
fails on knots too, where metabelian targets are easy. Knot counting is
$\sharpP$-hard for $S_4$, the smallest group for which this can happen, and for
its double covers $\mathrm{GL}(2,3)$ and $2O$. It is also $\sharpP$-hard for
groups of Fitting height two: for $\mathrm{UT}_3(\F_4)\rtimes
C_3$, and for the Heisenberg groups $H(\F_p)\rtimes C_r$ ($p$ odd) whose
action has a ``hard pair'' of eigencharacters. For the second family we prove
the converse: without a hard pair, the count is polynomial for knot groups and
for closed 3-manifolds with bounded first homology. The $S_4$ proof rests on a
``bicycle law'' for transposition colourings of alternating diagrams with
tripled crossings, which turns the count into a 12-state Potts-type model. The
Heisenberg proofs read the count as the kernel size of a twisted triple cup
product, which Borromean surgeries steer freely enough to encode Potts
partition functions. For $\mathrm{UT}_3(\F_4)\rtimes C_3$ the reduction also
preserves approximation, so approximating the count is $\#\mathrm{BIS}$-hard.
\end{abstract}

\section{Results}

For a finite group $G$ and a link $L\subset S^3$, write
\[
H_L(G)=|\Hom(\pi_1(S^3\setminus L),G)|.
\]
Up to normalization it is the link invariant of the modular category
$\mathrm{Rep}(D(G))$. Rowell
\cite[\S3.3.5]{Row09} made the following conjecture.

\begin{quote}
\textbf{Conjecture 3.2.} Let $G$ be a finite group and $L$ a link.
(a) There exists an FPRAS for computing $H_L(G)$ for any group $G$.
(b) Suppose $G$ is solvable. Then there is a polynomial algorithm for exact
computation of $H_L(G)$.
\end{quote}

Kuperberg and Samperton proved that for every nonabelian simple $G$ and
nontrivial conjugacy class $C$, counting the homomorphisms from knot groups to
$G$ that send the meridian into $C$ is $\sharpP$-complete \cite{KS21}. They also wrote that they were ``much less
confident'' in part~(b) \cite[\S5.3]{KS18}. Krovi and Russell had shown exact
$\sharpP$-hardness on links for the alternating groups $A_m$, $m\ge5$, with
meridians in a fixed class \cite{KR15}. On the
positive side, $H_K(G)=|G|$ for nilpotent $G$, and metabelian targets reduce
to linear algebra over the Alexander module \cite[\S1]{KS21}.
Galindo~\cite{Gal26} proved that $|\Hom(\pi_1M,G)|$ is $\sharpP$-hard for
every nonabelian $G$ when $M$ ranges over closed 3-manifolds; his hard
instances are graph manifolds whose first Betti number grows with the input.
Samperton and Wei\ss~\cite{SW26} proved that class-two nilpotent targets are
polynomial when $|H^2(M;Z(G))|$ is bounded, and named $S_3$ on 3-manifold groups
as open. None of these addresses solvable targets on link or knot complements,
and we found no statement in either direction.

On links, part~(b) fails for a simple reason (Proposition~\ref{prop:links}): the
components of a link may carry different conjugacy classes, and the sum over
these choices encodes the Potts model, already for $S_3$. This uses unboundedly
many components; for links with a bounded number of components every
metabelian target gives a polynomial-time count. On knots, whose groups have
abelianization $\Z$, nilpotent targets give $H_K(G)=|G|$ and metabelian
targets reduce to linear algebra, so a counterexample must have derived length
at least three. Our main results are such counterexamples.

All reductions below are polynomial-time Turing reductions, and links are
given by diagrams. Recall that $\FP=\sharpP$ would mean that every function
in $\sharpP$ is computable in polynomial time; it implies $\mathrm{P}=
\mathrm{NP}$.

\begin{proposition}[links]\label{prop:links}
For the dihedral group $D_n$ of order $2n$ with $n\ge3$ odd (so $D_3=S_3$),
computing $H_L(D_n)$ from a link diagram is $\sharpP$-hard. For each fixed $C$
and each metabelian $G$, $H_L(G)$ is computable in polynomial time on links with
at most $C$ components.
\end{proposition}

\begin{theorem}[$S_4$]\label{thm:s4}
For $G=S_4$, for its double covers $G=\mathrm{GL}(2,3)$ and $G=2O$ (the binary
octahedral group), and for the count of maps sending the meridian to a
transposition, computing $H_K(G)$ from a knot diagram is $\sharpP$-hard.
Every group of order less than $24$, and $\mathrm{SL}(2,3)$, gives a
polynomial-time count, so $S_4$ is the smallest group with this property
unless $\FP=\sharpP$.
\end{theorem}

Let $\omega$ be a primitive cube root of unity in $\F_4$. For a field $\F_q$
let $H(\F_q)$ be the Heisenberg group of triples with
$(x,y,z)(u,v,t)=(x+u,\,y+v,\,z+t+xv)$. For $r\mid q-1$ and
$\alpha,\beta\in\F_q^\times$ with $\alpha^r=\beta^r=1$, let
\[
G_{\alpha,\beta}=H(\F_q)\rtimes C_r,\qquad h\cdot(x,y,z)=(\alpha x,\ \beta y,\
\alpha\beta z),
\]
where $h$ generates $C_r$. The group $\mathrm{UT}_3(\F_4)\rtimes C_3$ of order
$192$ is $G_{\omega^2,\omega^2}$ over $\F_4$; it has trivial centre and Fitting
height two. Write $S^3_r(K)$ for $r$-surgery on $K$; its first homology is
$\Z/r$.

\begin{theorem}[$\mathrm{UT}_3(\F_4)\rtimes C_3$]\label{thm:ut3}
Let $G=\mathrm{UT}_3(\F_4)\rtimes C_3$. For every knot $K$,
\[
H_K(G)=|\Hom(\pi_1S^3_3(K),G)|+63.
\]
Computing either number from a diagram of $K$ is $\sharpP$-hard.
\end{theorem}

\begin{theorem}[Heisenberg dichotomy]\label{thm:heis}
Let $p$ be an odd prime, $r\ge2$ a divisor of $p-1$, and $G=G_{\alpha,\beta}$
over $\F_p$. Call $(\alpha,\beta)$ a \emph{hard pair} if $\alpha\ne1$,
$\beta\ne1$ and $\alpha\beta\ne1$.
\begin{enumerate}
\item If $(\alpha,\beta)$ is not a hard pair, then $H_K(G)$ is computable in
polynomial time from a knot diagram, and for every $C$ the number
$|\Hom(\pi_1M,G)|$ is computable in polynomial time for closed orientable
3-manifolds $M$ with $|H_1(M;\Z)|\le C$, given by triangulations.
\item If $(\alpha,\beta)$ is a hard pair, then computing $H_K(G)$, and
computing $|\Hom(\pi_1S^3_r(K),G)|$, from a diagram of $K$ is $\sharpP$-hard.
\end{enumerate}
\end{theorem}

For odd $p$ the smallest groups with a hard pair are $G_{2,2}$ over $\F_5$ with
$r=4$ (order $500$) and $G_{2,2}$ over $\F_7$ with $r=3$ (order $1029$); there is
none with $r=2$. Part~(1) is a special case of a more general statement
(Remark~\ref{rem:K}).

\begin{corollary}\label{cor:rowell}
Unless $\FP=\sharpP$, Rowell's Conjecture 3.2(b) fails on links for the
dihedral groups of Proposition~\ref{prop:links}, and already on knots for each of
the groups shown $\sharpP$-hard in Theorems~\ref{thm:s4}--\ref{thm:heis}.
\end{corollary}

\noindent\emph{Proof.} All these groups are solvable, and a polynomial
algorithm for $H_L(G)$ would compute a $\sharpP$-hard function in
polynomial time. \qed

\smallskip
For $\mathrm{UT}_3(\F_4)\rtimes C_3$ the reduction also preserves
approximation, so approximating $H_K(G)$ is $\#\mathrm{BIS}$-hard
(Proposition~\ref{prop:bis}); this bears on part~(a).

\paragraph{The line, as tested.} On knots, metabelian targets are easy, and
the hard examples above have derived length three. Derived length three does
not suffice: $\mathrm{SL}(2,3)$ and the groups $G_{\alpha,\beta}$ with
$\alpha\beta=1$, $\alpha\ne1$, have derived length three and are polynomial.
In the Heisenberg family the obstruction to lifting a map is a twisted
\emph{triple} product of first cohomology classes. Without a hard pair every
unbounded part of it is linear or pairwise and Gauss sums count its zeros;
with a hard pair it is a genuinely three-way form whose alternating part can be
steered at will, and that freedom encodes the Potts model. The same picture
separates $\mathrm{SL}(2,3)$ (quadratic obstruction, bounded target) from
$\mathrm{UT}_3(\F_4)\rtimes C_3$ (cubic obstruction).

\section{Preliminaries}\label{sec:prelim}

\paragraph{Wirtinger colourings.} A homomorphism from a link group to $G$ is a
colouring of the arcs of a diagram by elements of $G$ such that at each
crossing, with over-arc coloured $g$ and incoming under-arc coloured $x$, the
outgoing under-arc is coloured $gxg^{-1}$ or $g^{-1}xg$ according to the sign
of the crossing. The signs of individual crossings matter in general, but not
when $g$ is an involution, as for the transposition colourings of
Section~\ref{sec:s4}. Reversing the convention at every crossing at once
replaces each colouring by its pointwise inverse, so counts over conjugacy
classes closed under inversion do not depend on it. All meridians of a knot are conjugate and generate the group, so a
homomorphism from a knot group sends every meridian into one conjugacy class.

\paragraph{Counting over a quotient.} If $G=N\rtimes Q$ with $N$ abelian, the
homomorphisms $\pi\to G$ lying over a fixed $\psi\colon\pi\to Q$ are the
$1$-cocycles $Z^1(\pi;N_\psi)$, where $N_\psi$ is $N$ with $\pi$ acting through
$\psi$. For a central extension $1\to A\to\tilde G\to G\to1$, a map
$\rho\colon\pi\to G$ lifts if and only if an obstruction class in
$H^2(\pi;A)$ vanishes, and the lifts then form a torsor under $H^1(\pi;A)$;
twisted versions of both statements hold for extensions by modules
\cite{MS05}.

\paragraph{Potts and Tutte.} For a multigraph $\Gamma$ with $n$ vertices, $m$
edges and $k$ components, the Potts partition function is
\[
Z_\Gamma(Q,v)=\sum_{\sigma\colon V\to[Q]}\prod_{uv\in E}\bigl(1+v[\sigma_u=\sigma_v]\bigr)
=\sum_{S\subseteq E}Q^{k(S)}v^{|S|}=Q^kv^{n-k}\,T_\Gamma\!\left(1+\tfrac Qv,\,1+v\right),
\]
where $k(S)$ counts the components of $(V,S)$ and $T_\Gamma$ is the Tutte
polynomial. Evaluating $T_\Gamma$ at a rational point $(x,y)$ is
$\sharpP$-hard unless $(x-1)(y-1)=1$ or $(x,y)$ is one of eight special
points \cite{JVW90}; on planar graphs it stays hard off the hyperbolas
$(x-1)(y-1)\in\{1,2\}$ and finitely many points \cite{Ver05}. In particular the
chromatic polynomial $P(H,12)$ of a planar graph is $\sharpP$-hard to compute.
The hardness proofs use interpolation, so they are Turing reductions.

\section{Links: the dihedral groups}\label{sec:links}

For a link $L$ with a component $K_i$, write $\mu_i$ for its meridian and
$\lambda_i$ for its Seifert longitude, and sum over
$\rho\in\Hom(\pi_1(S^3\setminus L),G)$.

\begin{lemma}[cables and local knots]\label{lem:cable}
\begin{enumerate}
\item Replacing $K_i$ by $k$ parallel copies, pushed off along the Seifert
framing, gives $H=\sum_\rho|C_G(\rho(\lambda_i))|^{k-1}$.
\item Tying a knot $J$ into $K_i$ gives $H=\sum_\rho N_J(\rho(\mu_i))$, where
$N_J(g)$ counts the homomorphisms from the group of $J$ that send its meridian
to $g$. If $G$ is metabelian, the new longitude has image $\rho(\lambda_i)$.
\end{enumerate}
\end{lemma}

\begin{proof}
(1) The new exterior is the old one with the cable space
$(D^2\setminus k\text{ points})\times S^1$ glued to the boundary torus of $K_i$,
the boundary circle of the disk going to $\mu_i$ and the circle factor to
$\lambda_i$. Its group is $F_k\times\Z$, so an extension of $\rho$ chooses
$a_1,\dots,a_k\in C_G(\rho(\lambda_i))$ with $a_1\cdots a_k=\rho(\mu_i)$; since
$\rho(\mu_i)$ commutes with $\rho(\lambda_i)$, there are $|C_G(\rho(\lambda_i))|^{k-1}$
choices. (2) The new exterior is the union of the two exteriors along a
meridional annulus, so maps are pairs that agree on the meridian. The new
longitude is $\lambda_i\lambda_J$, and $\lambda_J$ lies in the second commutator
subgroup, which a metabelian $G$ kills.
\end{proof}

\begin{proof}[Proof of Proposition~\ref{prop:links}]
Let $\Gamma=(V,E)$ be a multigraph with $n$ vertices and $m$ edges, and
$c=n+m$. Let $L_\Gamma$ be a link with an unknotted component $U_v$ for each
vertex and $K_e$ for each edge, with $\mathrm{lk}(K_e,U_u)=1$ and
$\mathrm{lk}(K_e,U_v)=-1$ for $e=uv$ and all other linking numbers $0$; the
closure of the pure braid $\prod_{e=uv}A_{u,e}A_{v,e}^{-1}$ is one. Choose a
prime $p\mid n$ and let $J$ be the torus knot $T(2,p)$. A homomorphism from its
group to $D_n$ is determined by the meridian when the meridian is a rotation,
and there are $p$ of them when it is a fixed reflection, so $N_J$ is $1$ on
rotations and $p$ on reflections. Let $L(j)$ be $L_\Gamma$ with $j$ copies of
$J$ tied into every component and then every edge component replaced by two
parallel copies. By Lemma~\ref{lem:cable},
\[
H_{L(j)}(D_n)=\sum_{s=0}^{c}p^{js}B_s,\qquad
B_s=\sum_{\rho:\ s\text{ reflection meridians}}\ \prod_{e\in E}|C(\rho(\lambda_e))|.
\]
The $c+1$ values $j=0,\dots,c$ determine $B_0$ through a Vandermonde system in
the distinct nodes $p^s$. In $B_0$ every meridian goes to the cyclic subgroup
$C_n$ of rotations, so $\rho$ factors through $H_1=\Z^c$ and is any
$x\in(\Z/n)^c$; then $\rho(\lambda_e)=x_u-x_v$, whose centralizer has order $2n$
if $x_u=x_v$ and $n$ otherwise ($n$ is odd). Hence
\[
B_0=n^m\sum_{x\colon V\to\Z/n}\ \prod_{uv\in E}n\bigl(1+[x_u=x_v]\bigr)=n^{2m}Z_\Gamma(n,1),
\]
and $Z_\Gamma(n,1)=n^kT_\Gamma(1+n,2)$ lies on $(x-1)(y-1)=n\ge3$, which is
$\sharpP$-hard \cite{JVW90} (also for planar $\Gamma$ \cite{Ver05}). The link $L(j)$
has polynomial size. Using $k$-fold instead of $2$-fold cables and
interpolating in $k$ recovers $Z_\Gamma(n,2^{k-1}-1)$ for several $k$, hence the
number of proper $n$-colourings, without appeal to the Tutte classification.

For the second statement, the maps to the abelian quotient $G/G'$ factor
through $H_1=\Z^c$, so there are at most $|G/G'|^C$ of them, and the lifts of each
through the abelian $G'$ solve a linear system (\texttt{BOUNDED-H1-20261025-001B}).
\end{proof}

\begin{remark}\label{rem:galindo}
More is true if one accepts Galindo's theorem~\cite{Gal26}, whose hard
instances are closed manifolds given by planar surgery diagrams with framing
$0$. Replacing every component of such a diagram by $k$ framed parallels, for
polynomially many $k$, determines by interpolation the number of $\rho$ that
send every framed longitude into the centre of $G$, since only those reach the
largest node $|G|^c$ (Lemma~\ref{lem:cable}(1) with framings). For $G$ with
trivial centre this is the count of the surgered closed manifold, so
$H_L(G)$ is $\sharpP$-hard on links for every nonabelian $G$ with trivial
centre. We have not checked the proofs of~\cite{Gal26}. For nilpotent targets
the cables see framed longitudes only modulo the centre, and we do not know the
complexity on links: for $D_4$ and $Q_8$ the count depends only on the linking
numbers, as $4^c\sum_S4^{-\operatorname{cutrank}(S)}$ over subsets $S$ of the
components, where $\operatorname{cutrank}(S)$ is the rank over $\F_2$ of the
linking matrix with rows in $S$ and columns outside $S$ (verified on 60 random
links each).
\end{remark}

\section{The symmetric group \texorpdfstring{$S_4$}{S4}}\label{sec:s4}

\subsection{Reduction to transposition colourings}

Let $T(K)$ and $R(K)$ count the homomorphisms $\pi_1(S^3\setminus K)\to S_4$
that send the meridian to a transposition and to a $3$-cycle respectively.

\begin{lemma}\label{lem:split}
$H_K(S_4)=4+R(K)+2T(K)$, and $R(K)$ is computable in polynomial time.
\end{lemma}

\begin{proof}
Sort the maps by the class of the meridian's image. The identity gives one
map. If the meridian goes to a double transposition, the image lies in the
abelian normal subgroup $V_4$, so the map factors through $H_1=\Z$: three
maps. Livingston \cite[Thm.~9]{Liv95} proved that over each Fox
$3$-colouring the lifts to $S_4$ with a transposition meridian and with a
$4$-cycle meridian are equally many, so the $4$-cycle class contributes
$T(K)$; for colourings with cyclic image both lifts lie in a dihedral group
of order 8 and are counted directly. Finally a map with a $3$-cycle meridian
lands in $A_4=V_4\rtimes C_3$, which is metabelian, so $R(K)$ is a sum of two
cocycle counts $|Z^1(\pi;\F_{4,\omega^{\pm1}})|$, computed by linear algebra.
\end{proof}

The repository re-proves Livingston's identity
(\texttt{CLOSED-H1-20261025-001B}). It remains to show that $T(K)$ is
$\sharpP$-hard.

\subsection{Tripled alternating diagrams}

Let $G=(V,E)$ be a connected plane multigraph with $n$ vertices and $m$ edges.
Let $D(G)$ be the alternating link diagram whose Tait graph is $G$: the black
regions are the vertices, and each edge is one crossing. Let $G^{(3)}$ repeat
every edge three times side by side, so that $D(G^{(3)})$ has a twist of three
crossings for each edge. For a spanning subgraph $H$ write $c(H)$ for its
number of components (isolated vertices included) and $B(H)$ for its bicycle
space over $\F_2$: the edge sets that are both cycles and cocycles. Let
$\Col(D)$ be the number of colourings of $D$ by the six transpositions of
$S_4$.

The transpositions of $S_4$ lie over the transpositions of $S_3=S_4/V_4$, two
over each, and the two commute. Label the transpositions of $S_3$ by $\Z/3$ so
that conjugation reads $a\triangleright b=2a-b$. A twist of $k$ crossings acts
on the colours at its top by the $k$-th power of the Hurwitz move
$\sigma(x,y)=(x\triangleright y,\,x)$.

\begin{lemma}\label{lem:twist}
On pairs of transpositions of $S_3$, $\sigma^3$ is the identity. On pairs of
transpositions of $S_4$ over colours $a,b$: if $a=b$ then
$\sigma^3(x,y)=(y,x)$; if $a\ne b$ then $\sigma^3$ keeps the $S_3$-colours in
place and changes each of the two lifts by an element of $V_4$ that depends
only on $(a,b)$.
\end{lemma}

\begin{proof}
$\sigma^3(a,b)=(3b-2a,\,b)=(a,b)$ modulo $3$. Over equal colours the two
transpositions commute, so $\sigma(x,y)=(y,x)$. The remaining case is a
direct check on the $6\times6$ table of conjugates.
\end{proof}

\begin{proposition}[the bicycle law]\label{prop:bicycle}
\[
\Col\bigl(D(G^{(3)})\bigr)=\sum_{\varphi\colon V\to\Z/3}2^{\,c(H_\varphi)+\dim
B(H_\varphi)},\qquad H_\varphi=\bigl(V,\{uv\in E:\varphi_u=\varphi_v\}\bigr).
\]
Equivalently, with $\nu(t,\varphi,x)=(1-t,\varphi,1-x)$ the fixed-point-free
involution of $\Z/2\times\Z/3\times\Z/2\cong\Z/12$, $N$ its permutation matrix
and $J$ the all-ones matrix,
\[
\Col\bigl(D(G^{(3)})\bigr)=2^{-n}\sum_{z\in(\Z/12)^V}\ \prod_{uv\in
E}M(z_u,z_v),\qquad M=J-2N .
\]
\end{proposition}

\begin{proof}
By Lemma~\ref{lem:twist} each twist passes $S_3$-colours along positions, as
the smoothing that separates the two black regions would. That smoothing
turns $D(G^{(3)})$ into one circle around each black region, so the
$S_3$-colourings correspond to potentials $\varphi\colon V\to\Z/3$, and the
twist of $uv$ is monochromatic exactly when $\varphi_u=\varphi_v$.

Fix $\varphi$. A lift to $S_4$ chooses one of two transpositions on each arc.
By Lemma~\ref{lem:twist} the choices follow the strands through monochromatic
twists, which cross, and keep their positions, up to fixed elements of $V_4$,
through the other twists, which is the smoothing that deletes the edge. So the
lifts form an affine space whose dimension is the number of closed curves of
the diagram of $H_\varphi$, each isolated vertex contributing its circle. The
space is not empty: $S_3=\mathrm{Stab}(4)\subset S_4$ is a complement of $V_4$
whose transpositions are transpositions of $S_4$. By Rosenstiehl and Read
\cite{RR78}, \cite[\S17.3]{GR01}, a connected plane graph $C$ has
$1+\dim B(C)$ left-right cycles, so the number of curves is
$c(H_\varphi)+\dim B(H_\varphi)$.

For the second form, $c(H)+\dim B(H)$ is the nullity over $\F_2$ of the
Laplacian $L_H$: $x\mapsto\delta x$ maps $\ker L_H$ onto $B(H)$ with the
locally constant functions as kernel. The condition $L_Hx=0$ says that every
vertex has an even number of $H$-edges whose ends carry different bits.
Writing each parity condition as $\tfrac12\sum_{t_v\in\F_2}(-1)^{t_v(\cdots)}$
gives the edge factor $(-1)^{(t_u+t_v)[\varphi_u=\varphi_v][x_u\ne x_v]}
=1-2[z_v=\nu(z_u)]$ with $z=(t,\varphi,x)$.
\end{proof}

\begin{remark}\label{rem:links}
If every edge of $G$ is subdivided, summing out each middle vertex gives
$\tfrac12M^2=4J+2I$ per edge, and
$\Col=2^{-n}4^mZ_G(12,\tfrac12)$. The diagram is then a link with
$1+\beta_1(G)$ components, and one query to the colouring count evaluates the
planar Tutte polynomial at $(25,\tfrac32)$, which is $\sharpP$-hard. Knots need
the gadgets below.
\end{remark}

\subsection{Knots}

A two-terminal series-parallel network $\Gamma$, with every edge tripled,
substituted for each edge of $G$, gives an interaction matrix $M_\Gamma$ on
$\Z/12$ after its internal vertices are summed (with a factor $\tfrac12$ each).
Series composition is the matrix product (with the factor $\tfrac12$ for the
middle vertex) and parallel composition the entrywise product. The span of $J,I,N$ is closed under both, because $J^2=12J$,
$JN=NJ=J$, $N^2=I$, $J\circ X=X$, $N\circ N=N$ and $I\circ N=0$. Writing
$M_\Gamma=a\bigl((1+b)I+(1+g)N+(J-I-N)\bigr)$, parallel composition multiplies
$a$, $1+b$ and $1+g$.

\begin{lemma}\label{lem:knot}
\begin{enumerate}
\item If $\tau$ counts spanning trees, $t_\Gamma$ the spanning trees of $\Gamma$
and $f_\Gamma$ its spanning $2$-forests separating the terminals, then
$\tau(G[\Gamma])=\tau(G)\,t_\Gamma^{\,n-1}f_\Gamma^{\,m-n+1}$.
\item $D(K^{(3)})$ is a knot if and only if $\tau(K)$ is odd, for a connected
plane multigraph $K$.
\end{enumerate}
\end{lemma}

\begin{proof}
(1) A spanning tree of $G[\Gamma]$ meets each copy of $\Gamma$ in a spanning
tree or in a separating $2$-forest, and the copies of the first kind form a
spanning tree of $G$. (2) A twist with an odd number of crossings connects
strands like one crossing, so $D(K^{(3)})$ and $D(K)$ have the same number of
components, namely $1+\dim B(K)$ by Rosenstiehl--Read. And $\dim B(K)$ is the
$\F_2$-nullity of the reduced Laplacian, whose determinant is $\tau(K)$.
\end{proof}

The gadget $\Gamma_1=S(e,P(e,S(e,e)))$ has
$M_{\Gamma_1}=17J+8I-2N$ (with the factors $\tfrac12$ included), so
$(1+b,1+g)=(25/17,15/17)$, $t=3$ and $f=5$. Its $k$-fold parallel power
$\Gamma_1^{\parallel k}$ has $t_k=3k\cdot5^{k-1}$ and $f_k=5^k$, both odd for odd $k$.

\begin{proof}[Proof of Theorem~\ref{thm:s4}]
Let $H$ be a connected planar graph; $P(H,12)$ is $\sharpP$-hard. Double every
edge outside a spanning tree $T_0$ to get $\tilde H$, with $\tilde m$ edges.
Then $P(\tilde H,12)=P(H,12)$, and $\tau(\tilde H)=\sum_T2^{|T\setminus T_0|}$
is odd. For odd $k=1,3,\dots,2K-1$, with $K=(\tilde m+1)(\tilde m+2)/2$, the
diagram $D\bigl((\tilde H[\Gamma_1^{\parallel k}])^{(3)}\bigr)$ is a knot by
Lemma~\ref{lem:knot}, with $12k\tilde m$ crossings, and by
Proposition~\ref{prop:bicycle}
\[
T\bigl(D((\tilde H[\Gamma_1^{\parallel k}])^{(3)})\bigr)=
2^{-n}\,17^{k\tilde m}\sum_{i,j}N_{ij}\,\beta_{ij}^{\,k},
\qquad \beta_{ij}=\Bigl(\tfrac{25}{17}\Bigr)^{i}\Bigl(\tfrac{15}{17}\Bigr)^{j},
\]
where $N_{ij}$ counts the $z\in(\Z/12)^V$ with $i$ edges having $z_u=z_v$ and
$j$ edges having $z_v=\nu(z_u)$. The $\beta_{ij}=5^{2i+j}3^j17^{-i-j}$ are
pairwise distinct, so the $K$ counts determine every $N_{ij}$ through a
Vandermonde system in $\beta_{ij}^2$.

The reduction runs in polynomial time. It makes $K=O(\tilde m^2)$ queries,
and each diagram has $12k\tilde m=O(\tilde m^3)$ crossings and as many arcs,
so each count $T$ is at most $6^{12k\tilde m}$ and has $O(\tilde m^3)$ bits.
Each $\beta_{ij}$ is a quotient of integers with $O(\tilde m)$ bits and
$k\le2K-1=O(\tilde m^2)$, so every coefficient $\beta_{ij}^k$ and every factor
$2^{-n}17^{k\tilde m}$ has $O(\tilde m^3)$ bits. The coefficient matrix of the
system is a Vandermonde matrix in the $\beta_{ij}^2$ times the invertible
diagonal matrix of the $\beta_{ij}$, and exact Gaussian elimination over
$\mathbb{Q}$ solves it in time polynomial in these bit lengths. The reduction
is a polynomial-time Turing reduction with $K$ queries, not a parsimonious one.

Finally
$P(H,12)=\sum_jN_{0j}$: a colouring is proper exactly when no edge has equal
ends, whatever happens on the $\nu$-related edges. By Lemma~\ref{lem:split},
$T=(H_K(S_4)-4-R)/2$, so $H_K(S_4)$ is $\sharpP$-hard as well.

For the double covers $\tilde G\in\{\mathrm{GL}(2,3),2O\}$, central extensions
of $S_4$ by $\Z/2$: every map from a knot group to $S_4$ lifts, because
$H^2$ of a knot exterior vanishes, and the lifts form a torsor under
$H^1(\pi;\F_2)=\F_2$. So $H_K(\tilde G)=2H_K(S_4)$.

For the smallest-group statement: every group of order less than $24$ is
metabelian (it has an abelian normal Sylow subgroup with abelian quotient, or is
a $2$-group of order at most $16$, whose derived subgroup has order at most
$4$), and the only non-metabelian groups of order $24$ are $S_4$ and
$\mathrm{SL}(2,3)$. Metabelian (indeed abelian-by-nilpotent) targets give
polynomial counts on knots: maps to the nilpotent quotient are boundedly many
because $H_1=\Z$, and their lifts through the abelian kernel $A$ are the
solutions of a system of linear equations over $A$
(\texttt{BOUNDED-H1-20261025-001B}). Writing
$A=\Z/d_1\times\dots\times\Z/d_s$, these are integer congruences modulo the
$d_i$, and the Smith normal form of the integer matrix that combines the
coefficients with the moduli counts their solutions in polynomial time. This
covers primary parts with several different exponents, which need not be free
over a single $\Z/p^a$. For $\mathrm{SL}(2,3)$, a central
extension of $A_4$ by $\Z/2$, the argument for the double covers gives
$H_K(\mathrm{SL}(2,3))=2H_K(A_4)$.
\end{proof}

\section{Heisenberg groups twisted by a cyclic group}\label{sec:heis}

Throughout, $q=p$ is an odd prime with $r\mid p-1$, or $q=4$ with $r=3$, and
$G=G_{\alpha,\beta}$. Let $M$ be a closed oriented 3-manifold with
$\pi=\pi_1M$, and $\psi\colon\pi\to\Z/r$. For $\chi\in\mu_r(\F_q)$ let
$\F_\chi$ be $\F_q$ with $g\in\pi$ acting by $\chi^{\psi(g)}$, and
$V_\chi=H^1(M;\F_\chi)$. Put $\gamma=(\alpha\beta)^{-1}$.

\subsection{The count is the kernel size of a triple product}

\begin{proposition}[\texttt{COPRIME-CLASS2-20261026-001B}, \texttt{CORRECTIONS-20261026-001A}]\label{prop:count}
The homomorphisms $\pi\to G$ lying over $\psi$ number
\[
|Z^1(\pi;\F_{\alpha\beta})|\cdot|B^1(\pi;\F_\alpha)|\cdot|B^1(\pi;\F_\beta)|\cdot\cN_\psi,
\qquad \cN_\psi=\#\{(a,b)\in V_\alpha\times V_\beta:\ T(a,b,\cdot)=0\},
\]
where $T(a,b,c)=\langle a\smile b\smile c,[M]\rangle$ for $c\in V_\gamma$. For
$\mathrm{UT}_3(\F_4)\rtimes C_3$ and $H_1(M)=\Z/3$ this gives
$|\Hom(\pi,G)|=1+128\cdot4^d\cdot\cN_\psi$ with $d=\dim_{\F_4}V_{\omega^2}$.
\end{proposition}

\begin{proof}
The centre $A=\{(0,0,z)\}$ is invariant and $G/A=(\F_\alpha\oplus\F_\beta)
\rtimes C_r$, because $h$ preserves the section $(x,y)\mapsto(x,y,0)$. Maps to
$G/A$ over $\psi$ are pairs of cocycles $(c_x,c_y)$. The section has the
bilinear cocycle $((x,y),(u,v))\mapsto xv$, so the obstruction to lifting is
$[c_x\smile c_y]\in H^2(\pi;\F_{\alpha\beta})$, and when it vanishes the lifts
form a torsor under $Z^1(\pi;\F_{\alpha\beta})$. The class depends only on
the cohomology classes $a=[c_x]$, $b=[c_y]$, which accounts for the factors
$|B^1|$. Since $H^2(\pi)\to H^2(M)$ is injective and Poincar\'e duality pairs
$H^2(M;\F_{\alpha\beta})$ perfectly with $H^1(M;\F_\gamma)$, the class vanishes
exactly when $T(a,b,\cdot)=0$. For the last formula, $\psi=0$ gives only the
trivial map (a nontrivial image in the $2$-group $\mathrm{UT}_3(\F_4)$ would
need $2$-torsion in $H_1$), and the two surjections $\psi$ are Galois conjugate
and contribute $4\cdot4^d\cdot4^2\cdot\cN_\psi$ each.
\end{proof}

Equivalently $\cN_\psi=|V_\gamma|^{-1}\sum_{a,b,c}e(T(a,b,c))$ for a nontrivial
additive character $e$ of $\F_q$: the count is a cubic exponential sum.

\subsection{Without a hard pair: polynomial}

\begin{proof}[Proof of Theorem~\ref{thm:heis}(1)]
There are at most $|H_1(M)|$ maps $\psi$, and for each $k$ the characters
$\alpha^k,\beta^k,(\alpha\beta)^k$ still include $1$. The factors $|Z^1|$ and
$|B^1|$ are linear algebra. If $\alpha^k=1$ then $V_{\alpha^k}=H^1(M;\F_p)$ has
at most $|H_1(M)|$ elements: enumerate $a$ and count the kernel of the linear
map $b\mapsto T(a,b,\cdot)$. The case $\beta^k=1$ is symmetric. If
$(\alpha\beta)^k=1$, then $V_{\gamma^k}$ is bounded with a basis
$c_1,\dots,c_s$, and $\cN$ counts the common zeros of $s$ bilinear forms:
\[
\cN=p^{-s}\sum_{\theta\in\F_p^s}\ \sum_{a,b}e\Bigl(\textstyle\sum_i\theta_iT(a,b,c_i)\Bigr)
=p^{-s}\sum_{\theta}p^{\dim V_{\alpha^k}+\dim V_{\beta^k}-\operatorname{rank}T(\cdot,\cdot,\sum\theta_ic_i)}.
\]
All the inputs (cocycle spaces, cup products on a triangulation, the
obstruction modulo coboundaries) are polynomial-time linear algebra. For a
knot exterior $X$, which is aspherical, the count over $\psi$ is the same with
$\cN=\#\{(a,b):a\smile b=0\in H^2(X;\F_{(\alpha\beta)^k})\}$. Here $\psi$
ranges over the $r$ maps $\Z\to\Z/r$, $H^1(X;\F_p)=\F_p$ and $H^2(X;\F_p)=0$,
so each term is again an enumeration of a bounded space followed by linear
algebra.
\end{proof}

\begin{remark}\label{rem:K}
The same argument proves more (\texttt{COPRIME-CLASS2-20261026-001C}). Let
$G=N\rtimes C_r$ with $N$ a $p$-group, $p\nmid r$, having a $C_r$-invariant
central subgroup $A$ with $A$ and $N/A$ elementary abelian and with
$N$ described by an equivariant bilinear cocycle (hypotheses (K)). Over a
splitting field the commutator pairing has components
$\kappa_{\chi,\chi'}\colon\bar B_\chi\times\bar B_{\chi'}\to\bar
A_{\chi\chi'}$. If no component with $\chi$, $\chi'$ and $\chi\chi'$ all
nontrivial is nonzero, counting is polynomial with bounded first homology
(for $p=2$ the quadratic part is handled by Arf invariants). Whether the
converse holds for all such $N$ is open; Theorem~\ref{thm:heis} proves it for
Heisenberg groups. Groups with higher-exponent layers are not covered.
\end{remark}

\subsection{Steering the triple product by Borromean surgery}

A Y-graph (a clasper with three leaves) in $M$ determines a Borromean surgery
$M\rightsquigarrow M'$ \cite{Mat87,Hab00}. When the leaves $\ell_1,\ell_2,
\ell_3$ lie in $\ker\psi$, the cover and the cohomology are undisturbed.

\begin{proposition}[\texttt{COPRIME-CLASS2-20261026-001B}(c), \texttt{CORRECTIONS-20261026-001C}]\label{prop:steer}
If the leaves lie in $\ker\psi$, then $H_1(M')\cong H_1(M)$, $\psi$ extends,
every $V_\chi$ is unchanged, and for characters $\chi_1\chi_2\chi_3=1$
\[
T'(a,b,c)-T(a,b,c)=\varepsilon\sum_{\sigma\in S_3}\operatorname{sgn}(\sigma)\,
\lambda_{\sigma1}(a)\,\lambda_{\sigma2}(b)\,\lambda_{\sigma3}(c),
\]
where $\lambda_i(a)$ evaluates a cocycle representing $a$ on the based leaf
$\ell_i$, and $\varepsilon=\pm1$ is a fixed orientation convention.
\end{proposition}

\begin{proof}[Sketch]
Since $\psi$ vanishes on the boundary of a neighbourhood of the Y-graph, and
$H_1$ of that boundary surjects onto $H_1$ of both the neighbourhood and the
handlebody glued in, the $r$-fold cyclic cover is trivial over the surgery
region and the twisted cocycles extend; this gives the first statements.
Massuyeau's formula for the change of the triple cup product under a
Y-surgery \cite[Lemma~3.2]{Mas03}, applied to the $r$ lifts in the cover,
adds $r$ equal determinants, because the characters multiply to $1$; the
transfer divides by $r$, which is invertible in $\F_q$. Changing the basing
path multiplies the three leaf functionals by $\chi_1(g),\chi_2(g),\chi_3(g)$,
whose product is $1$.
\end{proof}

This is the finite-cover, $\F_q$ analogue of Moussard's formula for null
Y-graphs and equivariant triple intersections in the infinite cyclic cover
\cite[Prop.~2.5, Thm.~2.7]{Mou17}. Our leaves need only lie in $\ker\psi$.

\begin{lemma}[pure leaves; \texttt{COPRIME-CLASS2-20261026-001B}(d)]\label{lem:pure}
Let $q=p$ and suppose $\psi$ is onto. For each $\zeta\in\mu_r(\F_p)$ and each
functional $f$ on $V_\zeta$ there is a loop in $\ker\psi$ whose functional is
$f$ on $V_\zeta$ and $0$ on $V_\chi$ for every $\chi\ne\zeta$.
\end{lemma}

\begin{proof}
$H_1$ of the cover with $\F_p$ coefficients surjects onto
$\bigoplus_\chi V_\chi^*$, and every class comes from a loop in $\ker\psi$.
Pick $t\in\pi$ with $\psi(t)=1$ and apply the eigenprojector
$P_\zeta=r^{-1}\sum_j\zeta^{-j}t^j$, realized by the product of the conjugates
$t^j\ell t^{-j}$ raised to the powers $r^{-1}\zeta^{-j}\bmod p$.
\end{proof}

The hypothesis that $\psi$ is onto is needed: for the lens space $L(7,1)$ with
$p=7$, $r=3$ and $\psi=0$ the three local systems coincide and no loop
separates them. In the constructions below $\psi$ sends the meridian to $1$.

\subsection{Gadgets}

Fix a multigraph $\Gamma$ with $n$ vertices and $m$ edges, and orient each
edge $e=uv$. Use coordinates $x_v$ ($v\in V$) and $p_e,s_e,w_e$ ($e\in E$),
and put $\ell_e=x_u-x_v$. For three functionals write $\lambda\wedge\mu\wedge\nu$
for the trilinear form $\det(\lambda,\mu,\nu)$, and for two
$(\lambda\wedge\mu)(a,b)=\lambda(a)\mu(b)-\mu(a)\lambda(b)$. Proposition~\ref{prop:steer}
can change the component $(\alpha,\beta,\gamma)$ only by forms that are
alternating in every pair of slots with equal characters. The coincidences
among $\alpha,\beta,\gamma$ give four patterns.

\begin{lemma}[\texttt{COPRIME-CLASS2-20261026-001D} \S1, \texttt{CORRECTIONS-20261026-001D} Lemma~6]\label{lem:gadget}
\begin{center}
\begin{tabular}{@{}llll@{}}
\toprule
Pattern & Spaces & Form $T$ & $\cN$\\
\midrule
$\alpha,\beta,\gamma$ distinct & $X\oplus P\oplus S$ in each slot & $\sum_e\ell_e\wedge p_e\wedge s_e$ & $q^{2m-n}Z_\Gamma(q^3,q^2-1)$\\
$\alpha=\beta=\gamma$ & $V_\alpha=X\oplus P\oplus S$ & $\sum_e\ell_e\wedge p_e\wedge s_e$ & $q^{2m-n}Z_\Gamma(q^3,q^2-1)$\\
$\alpha=\beta\ne\gamma$ & $V_\alpha=X\oplus P$, $V_\gamma=W$ & $\sum_e(\ell_e\wedge p_e)\otimes w_e$ & $q^{m}Z_\Gamma(q^2,q-1)$\\
$\gamma=\alpha\ne\beta$ & $V_\alpha=X\oplus P$, $V_\beta=S$ & $\sum_e s_e(b)\,(\ell_e\wedge p_e)(a,c)$ & $q^{m-n}Z_\Gamma(q^2,q-1)$\\
\bottomrule
\end{tabular}
\end{center}
For $q=4$, where a hard pair forces $\alpha=\beta=\gamma$, the second row holds
with $\ell_e=x_u+x_v$. The case $\gamma=\beta\ne\alpha$ is symmetric to the
last.
\end{lemma}

\begin{proof}
Write $\cN=|V_\gamma|^{-1}\sum_{a,b,c}e(T(a,b,c))$ and sum the edge coordinates
first; the sums factor over the edges. For the three-way form fix the vertex
parts and let $u_e=(\ell_e(a),\ell_e(b),\ell_e(c))\in\F_q^3$. Then
$\sum_{P\in\F_q^{3\times2}}e(\det(u_e\mid P))$ is $q^6$ if $u_e=0$ and $q^4$
otherwise: move $u_e$ to the first basis vector by $\mathrm{GL}_3$, after which
$\det=p_2s_3-p_3s_2$ and each product sum gives $q$. Colouring each vertex by
$(a_v,b_v,c_v)\in\F_q^3$, the total is
$q^{4m}\sum_\sigma\prod_e(1+(q^2-1)[\sigma_u=\sigma_v])$, and dividing by
$|V|=q^{n+2m}$ gives the first two rows. For $\alpha=\beta\ne\gamma$ the
condition at $e$ is $\ell_e(a)p_e(b)-p_e(a)\ell_e(b)=0$, which has $q^2$
solutions $(p_e(a),p_e(b))$ if $\ell_e(a)=\ell_e(b)=0$ and $q$ otherwise. For
$\gamma=\alpha\ne\beta$ the conditions are $s_e(b)\ell_e(a)=0$ for each edge
and $\sum_{e\ni w}\pm s_e(b)p_e(a)=0$ for each vertex; with $S=\operatorname{supp}s$
the second system has $q^{m-n+k(S)}$ solutions $p(a)$, and the first forces
$S$-edges to be monochromatic, so $\cN=q^{m-n}\sum_{S}(q-1)^{|S|}q^{2k(S)}$.
\end{proof}

All four Tutte points $(1+q^3/(q^2-1),\,q^2)$ and $(1+q^2/(q-1),\,q)$ lie on
$(x-1)(y-1)=q^3$ or $q^2$, so evaluating the Potts functions above is
$\sharpP$-hard \cite{JVW90}.

\subsection{Realizing the gadgets by knots}

\paragraph{Base knots.} For odd $p$ and $\zeta\in\F_p^\times\setminus\{1\}$,
the twist knot $K_m$ with Alexander polynomial $mt^2-(2m+1)t+m$ and
$m\equiv\zeta(\zeta-1)^{-2}\pmod p$ has $\zeta$ and $\zeta^{-1}$ as roots
modulo $p$. Its Alexander module is cyclic, so for $S^3_r(K_m)$ with $\psi$
the meridian character, $V_\zeta$ and $V_{\zeta^{-1}}$ are one-dimensional and
all other $V_\chi$ vanish. (For example $6_1$ has roots $2$ and $1/2$ for
every $p\ge5$.) For $q=4$ use the figure-eight knot, with $V_{\omega}$ and
$V_{\omega^2}$ one-dimensional over $\F_4$. In a connected sum the twisted
cohomology and $T$ split over the summands, because the splitting annulus has
no twisted cohomology. Every component of $T$ vanishes on the base: a nonzero
component would need three characters with product $1$ from one summand's
$\{\zeta,\zeta^{-1}\}$, forcing $\zeta^3=1$ on a one-dimensional space, where an
alternating form is zero; for $4_1$ over $\F_4$ the form is zero by direct
computation. With a hard pair no two of $\alpha,\beta,\gamma$ are inverse to
each other, so the summands for different slots do not interfere.

\paragraph{Construction} (\texttt{COPRIME-CLASS2-20261026-001D} \S3,
\texttt{CORRECTIONS-20261026-001D} \S3). Given $\Gamma$ and the pattern:
\begin{enumerate}
\item Take one base summand for each coordinate of each \emph{distinct}
character space of the pattern, drawn in disjoint balls $B_i$ joined by bands,
with a base point and a reference path into each ball.
\item In each $B_i$ pick arcs $a_i,b_i$ on which the summand's nonzero class
takes different values (constant cochains are coboundaries), and let
$\gamma_i$ be the loop around $a_i$ positively and $b_i$ negatively. It bounds
a disk in $B_i$ meeting the knot twice, and its functional is nonzero on the
summand's block. For odd $p$ apply Lemma~\ref{lem:pure} inside $B_i$ to get a
\emph{pure} leaf supported on the slot's eigenspace.
\item Realize $\ell_e$ as a band sum of the leaves of the summands of $u$ and
$v$ (the second reversed), choosing the band so that the closed loop through
it has $\psi\equiv0$; otherwise unbalanced gains change the count.
\item Expand the gadget into decomposable terms and give each term one
Y-graph whose leaves are parallel push-offs of the leaves it names, with edges
routed so that every leaf, read from the Y-graph's vertex, has its reference
value, and every leaf has the framing of its disk.
\item Each term has a leaf carrying a coordinate that occurs in no other
factor of that term ($p_e$, $s_e$ or $w_e$). Choose its cap disk disjoint from
all other Y-graphs; parallel caps suffice when another term uses the same
coordinate. Forgetting the knot, that leaf bounds a disk, so the Y-surgery
gives back $S^3$ \cite[Props.~3.3--3.4]{Hab00}, and the diffeomorphism is the
identity near a meridian. Hence the result is $S^3_r(K_\Gamma)$ for a knot
$K_\Gamma$, with the same framing.
\end{enumerate}
The Y-graphs and their leaf disks are pairwise disjoint, so the surgeries are
independent tangle replacements, and $K_\Gamma$ has $O(m(n+m))$ crossings.
By Proposition~\ref{prop:steer}, applied once per Y-graph, the component
$(\alpha,\beta,\gamma)$ of $T$ on $S^3_r(K_\Gamma)$ is the gadget form.

\paragraph{The other characters.} By Proposition~\ref{prop:count},
$|\Hom(\pi_1S^3_r(K_\Gamma),G)|=1+\sum_{k=1}^{r-1}c_k\cN_k$, where $c_k>0$
depends only on dimensions and $\cN_k$ reads the component
$(\alpha^k,\beta^k,\gamma^k)$. With pure leaves only components that are
permutations of $(\alpha,\beta,\gamma)$ change, and each such $\cN_k$ equals the
gadget count; every other component keeps its base value $0$, so its
$\cN_k=|V_{\alpha^k}||V_{\beta^k}|$. Over $\F_4$ the two surjections are Galois
conjugate and contribute equally. In all cases the oracle answer is an affine
function of the gadget count with known coefficients.

\subsection{Knot groups}

\begin{proof}[Proof of Theorem~\ref{thm:ut3}]
Let $N=\mathrm{UT}_3(\F_4)$ and $\rho\colon\pi_1(S^3\setminus K)\to G$, with
meridian $\mu$ and preferred longitude $\lambda$. If $\rho(\mu)\in N$, the
image is a finite $2$-group normally generated by one element, hence cyclic
(its abelianization is cyclic), so there are $|N|=64$ such maps; for them
$\rho(\lambda)=1$, and $\rho(\mu)^3=1$ forces $\rho(\mu)=1$, so only the trivial
one satisfies $\mu^3\lambda=1$. Otherwise $\lambda$ lies in the second
commutator subgroup, so $\rho(\lambda)\in G''=Z(N)$; it commutes with
$\rho(\mu)\in Nh^{\pm1}$, which acts on $Z(N)$ by $\omega^{\pm1}$ without fixed
points, so $\rho(\lambda)=1$. Every element of $Nh^{\pm1}$ has order $3$, so
$\rho(\mu)^3\rho(\lambda)=1$ holds automatically. Hence
$H_K(G)=|\Hom(\pi_1S^3_3(K),G)|+64-1$. Hardness: given $\Gamma$, build
$K_\Gamma$ with base $\#^{\,n+2m}4_1$, query once, and solve
$1+128\cdot4^{n+2m}\cdot4^{2m-n}Z_\Gamma(64,15)$ for $Z_\Gamma(64,15)$, a
$\sharpP$-hard Tutte evaluation at $(79/15,16)$.
\end{proof}

\begin{proof}[Proof of Theorem~\ref{thm:heis}(2)]
For the surgeries this is the construction above with one oracle call per
Tutte evaluation. For the knot exterior $X$ of $K_\Gamma$: if a character $\chi$ is
nontrivial on the meridian, the boundary torus and the surgery solid torus have
no $\chi$-twisted cohomology, so $H^i(X;\F_\chi)=H^i(S^3_r(K_\Gamma);\F_\chi)$
for $i=1,2$, compatibly with cup products. Hence a term of the count whose three
characters are all nontrivial on the meridian is the same for $X$ as for
$S^3_r(K_\Gamma)$. Terms in which one of
the three characters is trivial fall under the algorithm of part~(1). So
$H_{K_\Gamma}(G)$ is a computable quantity plus a known positive multiple of
the gadget count.
\end{proof}

\section{Remarks and open questions}\label{sec:remarks}

\paragraph{What is tested.} For links, $B_0=9^mZ_\Gamma(3,1)$ is recovered from
directly computed counts $H_{L(j)}(S_3)$ on five graphs (the largest link has
$2088$ crossings and $12$ components), the proper $3$-colourings of a path and a
triangle ($12$ and $6$) are recovered from link counts, and the same holds for
$D_5$ and $D_7$ on three graphs; the counting method is checked against brute
force on random braids, and Lemma~\ref{lem:cable} on separate cases. The
identities of Sections~\ref{sec:s4} and \ref{sec:heis} are checked by brute
force on small cases (Section~\ref{sec:evidence}).
The steering law is checked in all four patterns against direct computation of
the obstruction on braid closures with Borromean insertions. Complete
graph-to-knot instances were computed for the triangle and the digon over
$\F_4$ ($\cN=45162496$ and $326656$), for patterns two, three and four over
$\F_5$ and $\F_7$ ($\cN=3240625$, $46101601$ and $25925$), with leaf
configurations found by search and equal to the gadget up to a change of
basis. The band-and-projector recipe of the construction has not been run end
to end, and pattern one (which needs $r\ge6$) has steering checks but no
complete instance.

\paragraph{Approximate counting.} Part~(a) of Rowell's conjecture asks for an
FPRAS for every $G$. All our reductions land on ferromagnetic Potts points
($v>0$), and for one group the reduction preserves approximation.

\begin{proposition}\label{prop:bis}
For $G=\mathrm{UT}_3(\F_4)\rtimes C_3$, approximating $H_K(G)$ is
$\#\mathrm{BIS}$-hard under approximation-preserving reductions. So part~(a)
of Rowell's conjecture for this $G$ would give an FPRAS for $\#\mathrm{BIS}$.
\end{proposition}

\begin{proof}
By the proofs of Theorem~\ref{thm:ut3} and Proposition~\ref{prop:count}, the knot
$K_\Gamma$ built from a multigraph $\Gamma$ with $m$ edges has
$H_{K_\Gamma}(G)=64+128\cdot4^{4m}Z_\Gamma(64,15)$: over $\F_4$ both nontrivial
characters are steered, so no unsteered term remains. One query with relative
error $\delta$, minus the known constant $64$ and divided by $128\cdot4^{4m}$,
approximates $Z_\Gamma(64,15)$ with relative error at most $2\delta$. Approximating
the ferromagnetic Potts partition function $Z_\Gamma(Q,v)$ with $Q>2$ and
$v>0$ is $\#\mathrm{BIS}$-hard on general graphs \cite[Thm.~1]{GJ12}.
\end{proof}

This is evidence against part~(a), not a refutation. The other reductions do
not preserve approximation as written. For the Heisenberg groups with pure
leaves, the component read by $\psi^{-1}$ is never a permutation of
$(\alpha,\beta,\gamma)$, so it keeps its maximal value and dominates the count.
For $S_4$ the knot reduction interpolates, and the link reduction of
Remark~\ref{rem:links} lands on the planar ferromagnetic Potts model, for which
we know no $\#\mathrm{BIS}$-hardness. The interactions induced by kernel counts
are positive semidefinite, which suggests that this route gives at most
$\#\mathrm{BIS}$-type hardness of approximation.

\paragraph{Open questions.}
\begin{enumerate}
\item The converse of Remark~\ref{rem:K} for all coprime class-two groups with
a hard pair, and groups with higher-exponent layers.
\item Extension fields other than $\F_4$; for $q=2^k$ in patterns 2 and 4 a
twisted Bockstein term can appear.
\item Approximation for the Heisenberg groups and for $S_4$; upper bounds
($H_K(G)\le_{\mathrm{AP}}\#\mathrm{BIS}$) for any solvable $G$.
\item Closed 3-manifolds with bounded first homology for $\mathrm{GL}(2,3)$
and $2O$ (for $S_4$ itself the repository proves hardness with $H_1=\Z/2$,
\texttt{CLOSED-H1-20261025-001C}).
\end{enumerate}

\section{Evidence and replay}\label{sec:evidence}

The proofs, each with a standard-library check, are in the public
repository \url{https://github.com/jscul/knotlab} under
\texttt{proofs/<ticket>/}. The replay driver \texttt{replay.py} in
\texttt{proofs/ROWELL-NOTE-20261026-001/} copies them to a temporary directory,
reruns every check there, and compares each fresh result with the frozen one;
the frozen files are not touched.

\begin{center}
\begin{tabular}{@{}p{0.43\textwidth}p{0.5\textwidth}@{}}
\toprule
Claim & Proof and check\\
\midrule
Links: $S_3$ and $D_n$ ($n$ odd) hard; cables and local knots & \texttt{notes/knot-group-counting/links/} (\texttt{run\_all.py})\\
Bicycle law, Potts$(12,\tfrac12)$, knot gadgets, $T(K)$ $\sharpP$-hard & \texttt{BOUNDED-H1-20261025-001A}\\
Abelian-by-nilpotent targets polynomial on knots & \texttt{BOUNDED-H1-20261025-001B}\\
Livingston's identity re-proved & \texttt{CLOSED-H1-20261025-001B}\\
$H_K(S_4)=4+R+2T$; double covers & \texttt{SL23-BOUNDED-20261025-001B}\\
Count as a triple-product kernel ($\F_4$) & \texttt{CORRECTIONS-20261026-001A}\\
Steering ($\F_4$) & \texttt{CORRECTIONS-20261026-001C}\\
$\mathrm{UT}_3(\F_4)\rtimes C_3$ hard; knot groups ($+63$) & \texttt{CORRECTIONS-20261026-001D}\\
Count, steering and pure leaves (odd $p$) & \texttt{COPRIME-CLASS2-20261026-001B}\\
No hard pair: polynomial; Remark~\ref{rem:K} & \texttt{COPRIME-CLASS2-20261026-001C}\\
Hard pair: gadgets, construction, knot groups & \texttt{COPRIME-CLASS2-20261026-001D}\\
\bottomrule
\end{tabular}
\end{center}

An independent review with fresh code (\texttt{COPRIME-CLASS2-20261026-001A})
reproduced every number of the $\F_4$ case and 894 further steering
configurations; a second review replayed the four gadget patterns and the
$+63$ identity with separate code. Both were performed by AI systems and are
no substitute for expert review.

\begin{thebibliography}{99}
\bibitem{Gal26} C.~Galindo, \emph{A complexity dichotomy for quantum invariants of
3-manifolds}, arXiv:2605.06984 (2026).
\bibitem{GJ12} L.~A.~Goldberg and M.~Jerrum, \emph{Approximating the partition
function of the ferromagnetic Potts model}, J. ACM 59 (2012), art.~25.
\bibitem{GR01} C.~Godsil and G.~Royle, \emph{Algebraic Graph Theory}, Springer,
2001.
\bibitem{Hab00} K.~Habiro, \emph{Claspers and finite type invariants of links},
Geom. Topol. 4 (2000), 1--83.
\bibitem{JVW90} F.~Jaeger, D.~L.~Vertigan and D.~J.~A.~Welsh, \emph{On the
computational complexity of the Jones and Tutte polynomials}, Math. Proc.
Cambridge Philos. Soc. 108 (1990), 35--53.
\bibitem{KR15} H.~Krovi and A.~Russell, \emph{Quantum Fourier transforms and the
complexity of link invariants for quantum doubles of finite groups}, Comm. Math.
Phys. 334 (2015), 743--777.
\bibitem{KS18} G.~Kuperberg and E.~Samperton, \emph{Computational complexity and
3-manifolds and zombies}, Geom. Topol. 22 (2018), 3623--3670.
\bibitem{KS21} G.~Kuperberg and E.~Samperton, \emph{Coloring invariants of knots
and links are often intractable}, Algebr. Geom. Topol. 21 (2021), 1479--1510.
\bibitem{Liv95} C.~Livingston, \emph{Lifting representations of knot groups}, J.
Knot Theory Ramifications 4 (1995), 225--234.
\bibitem{Mas03} G.~Massuyeau, \emph{Cohomology rings, Rochlin function, linking
pairing and the Goussarov--Habiro theory of three-manifolds}, Algebr. Geom.
Topol. 3 (2003), 1139--1166.
\bibitem{Mat87} S.~V.~Matveev, \emph{Generalized surgery of three-dimensional
manifolds and representations of homology spheres}, Math. Notes 42 (1987),
651--656.
\bibitem{Mou17} D.~Moussard, \emph{Equivariant triple intersections}, Ann. Fac.
Sci. Toulouse Math. (6) 26 (2017), 601--643.
\bibitem{MS05} D.~Matei and A.~I.~Suciu, \emph{Counting homomorphisms onto finite
solvable groups}, J. Algebra 286 (2005), 161--186.
\bibitem{Row09} E.~C.~Rowell, \emph{Two paradigms for topological quantum
computation}, in: Advances in Quantum Computation, Contemp. Math. 482, Amer.
Math. Soc., 2009, 165--177.
\bibitem{RR78} P.~Rosenstiehl and R.~C.~Read, \emph{On the principal edge
tripartition of a graph}, Ann. Discrete Math. 3 (1978), 195--226.
\bibitem{SW26} E.~Samperton and A.~Wei\ss, \emph{Obstruction theory and the
complexity of counting group homomorphisms}, arXiv:2602.02885 (2026).
\bibitem{Ver05} D.~Vertigan, \emph{The computational complexity of Tutte
invariants for planar graphs}, SIAM J. Comput. 35 (2005), 690--712.
\end{thebibliography}

\end{document}
